% ============================================================
%  main.tex — Methodology only (Ch.4 extract for Overleaf, etc.)
%  Compile: pdfLaTeX or XeLaTeX
% ============================================================
\documentclass[11pt,a4paper]{article}

% ---- core packages ----
\usepackage[margin=1in]{geometry}
\usepackage{graphicx}
\usepackage{amsmath,amssymb,amsthm,bm}
\usepackage{booktabs}
\usepackage{multirow}
\usepackage{array}
\usepackage{siunitx}
\usepackage[hidelinks]{hyperref}
\usepackage{cite}
\usepackage{caption}
\usepackage{subcaption}
\usepackage{algorithm}
\usepackage{algpseudocode}
\usepackage{xcolor}
\usepackage{microtype}

% ---- figure path (place PDFs/PNGs next to this .tex, e.g. on Overleaf) ----
\graphicspath{{./}}

% ---- custom macros ----
\newcommand{\Pof}{\mathbf{P}^{\mathrm{OOF}}}
\newcommand{\logit}{\operatorname{logit}}
\newcommand{\RR}{\mathbb{R}}
\newcommand{\EE}{\mathbb{E}}
\newcommand{\ptilde}{\tilde p}
\newcommand{\phat}{\hat p}
\newcommand{\taustar}{\tau^{\star}}
\DeclareMathOperator*{\argmin}{arg\,min}
\DeclareMathOperator*{\argmax}{arg\,max}

% ---- title block ----
\title{An Actionable Attrition Prediction System:\\
Noise-Robust Stacking and Hard-Constrained Counterfactual\\
Interventions for New-Graduate Employees}
\author{Yuhang Song \and Collaborators \\ \small{\today}}
\date{}

% ============================================================
\begin{document}
\maketitle

% ============================================================
\section{Methodology}
\label{sec:methods}

% ------------------------------------------------------------
\subsection{Overview}
\label{sec:methods-overview}

We view the system as two tightly coupled artifacts: (i) an
\emph{inference pipeline} that maps a raw employee record to a binary
risk alert, and (ii) a \emph{training protocol} that estimates every
parameter consumed at inference time. Figure~\ref{fig:inference}
illustrates the former and Figure~\ref{fig:training} the latter. The two
figures are linked by a shared feature space $\varphi(\cdot)$, a shared
sample-weight vector $\mathbf{w}$, and a shared bundle of frozen
artifacts $\{\text{base}_k,\bm{\beta},b,g,\taustar\}$; every object in the
bundle has exactly one consumption site in the inference pipeline and
exactly one production site in the training protocol.

% ------------------------------------------------------------
\subsection{Feature Transform \texorpdfstring{$\varphi$}{phi}}
\label{sec:methods-phi}

Let $\mathbf{x}\in\RR^{12}$ denote the raw input of twelve integer-coded
survey variables (three demographic, four job-characteristic, five
work-environment). We apply a deterministic
transform $\varphi:\RR^{12}\to\RR^{41}$ that concatenates six
complementary subsets: (1) twelve cumulative Likert dummies
$\mathbf{1}\{x_j\ge c\}$ for $c\in\{2,3,4,5\}$ on the ordinal items;
(2) the twelve raw numerics; (3) six hand-crafted binary interactions
(e.g.\ income $\times$ pressure, satisfaction $\times$ match);
(4) eight expectation--realization gap features;
(5) two income transforms $\log(\text{income})$ and $\text{income}^2$;
and (6) a single KFold-target-encoded composite score. Tree-based learners
consume $\varphi(\mathbf{x})$ directly; distance- and gradient-based
learners additionally undergo z-score standardization
$\tilde{\varphi}(\mathbf{x}) = (\varphi(\mathbf{x}) - \bm{\mu}_{\mathrm{train}})/
\bm{\sigma}_{\mathrm{train}}$, where
$\bm{\mu}_{\mathrm{train}},\bm{\sigma}_{\mathrm{train}}$ are estimated
exclusively on training-fold samples to prevent leakage.

% ------------------------------------------------------------
\subsection{Inference Pipeline}
\label{sec:methods-inference}

The inference pipeline (Figure~\ref{fig:inference}) proceeds in six stages.
Given an input $\mathbf{x}$, each of the five base learners
$k\in\{\mathrm{RF}, \mathrm{NRB}, \mathrm{MT}, \mathrm{SVM}, \mathrm{ET}\}$
emits a probability $p_k = f_k(\varphi(\mathbf{x}))\in[0,1]$. A Level-1
L2-regularized logistic regression meta-learner aggregates the five base
probabilities in logit-space:
\begin{equation}
\label{eq:meta}
\ptilde \;=\; \sigma\!\left(\sum_{k=1}^{5}\beta_k\,\logit(p_k) + b\right),
\qquad
\sigma(z)=\frac{1}{1+e^{-z}}.
\end{equation}
The meta output $\ptilde$ is passed through a post-hoc isotonic calibration
map $g:[0,1]\to[0,1]$ (a monotone non-decreasing step function learned
via the Pool-Adjacent-Violators algorithm) to obtain the calibrated risk
score
\[
\phat \;=\; g(\ptilde).
\]
A final balanced-accuracy-optimal threshold $\taustar$ converts $\phat$
into a binary alert:
\[
\hat y \;=\; \mathbf{1}\{\phat \ge \taustar\}, \qquad \taustar = 0.135.
\]
All quantities on the right-hand sides are frozen at deployment;
Section~\ref{sec:methods-training} describes how each is estimated.

\begin{figure}[t]
  \centering
  \includegraphics[width=\linewidth]{fig_A_inference.pdf}
  \caption{\textbf{Inference pipeline of the proposed stacking ensemble.}
  A raw input $\mathbf{x}\in\RR^{12}$ is transformed by the deterministic
  map $\varphi:\RR^{12}\!\to\!\RR^{41}$ and distributed along two
  branches: tree-based learners (RF, NR-Boost, ET) consume
  $\varphi(\mathbf{x})$ directly, whereas distance/gradient-based learners
  (SVM-RBF, MT-MLP) consume its z-score--standardized version. Each base
  learner $k$ emits a probability $p_k\in[0,1]$ (colored circles). An
  L2-regularized logistic regression meta-learner aggregates the five
  probabilities in logit-space (Equation~\ref{eq:meta}); edge colors
  encode the sign of $\beta_k$ (blue $>\!0$, red $<\!0$) and edge widths
  encode $|\beta_k|$. The meta output $\ptilde$ is mapped to a calibrated
  risk score $\phat$ by an isotonic regressor $g$, and a
  balanced-accuracy-optimal threshold $\taustar=0.135$ produces the
  binary alert $\hat y$. Tensor shapes are annotated on every arrow;
  circled values illustrate the full trace of an exemplar high-risk
  employee.}
  \label{fig:inference}
\end{figure}

% ------------------------------------------------------------
\subsection{Training Protocol}
\label{sec:methods-training}

The training protocol (Figure~\ref{fig:training}) specifies how each
deployment artifact is estimated and decomposes into four stages.

\paragraph{Stage 1: Cleanlab-v6 sample weighting.}
To mitigate label noise that is endemic to self-reported turnover
surveys, we compute a self-confidence score for each training sample
using a noise-robust probe classifier fit on an independent
cross-validation loop. Samples in the bottom 10\% (ranked ascending) are
flagged as potentially mis-labeled and assigned a down-weight
$w_i = 0.3$; the remaining 90\% retain $w_i = 1.0$. The resulting binary
weight vector $\mathbf{w}\in\RR^{N}$ with $N=4{,}375$ is cached on disk
and consumed by every downstream training step.

\paragraph{Stage 2: 5-fold Out-of-Fold (OOF) training.}
For each base learner $f_k$, we run 5-fold stratified cross-validation
with $\mathbf{w}$ injected via the native \texttt{sample\_weight} API.
For each fold, the learner is trained on the four hold-in folds and
evaluated on the held-out fold; the concatenation of all held-out
predictions across the five folds yields a leakage-free OOF probability
matrix
\[
\Pof \;=\; \bigl[\,p_k^{\mathrm{OOF}}(\mathbf{x}_i)\,\bigr]_{i,k}
\;\in\; [0,1]^{N\times 5}.
\]
This matrix is the training input for the meta-learner; by construction,
each entry $p_k^{\mathrm{OOF}}(\mathbf{x}_i)$ is produced by a base
learner that has not seen sample $i$ during its fit, eliminating the
Level-1 leakage that plagues naive stacking implementations.

\paragraph{Stage 3: Meta fit and isotonic calibration.}
We logit-transform each column of $\Pof$ to obtain meta-features
$\mathbf{l}_i = \logit(\mathbf{p}_i^{\mathrm{OOF}})\in\RR^{5}$ and fit
the L2-regularized logistic regression by minimizing the weighted
cross-entropy with L2 penalty:
\begin{equation}
\label{eq:meta-objective}
\{\hat{\bm\beta},\hat b\}
\;=\;
\argmin_{\bm\beta,\,b}\;
-\sum_{i=1}^{N} w_i
\bigl[\,y_i\ln\hat p_i + (1-y_i)\ln(1-\hat p_i)\,\bigr]
\;+\;\frac{1}{2C}\|\bm\beta\|_2^{2},
\end{equation}
where $\hat p_i = \sigma(\bm\beta^{\!\top}\mathbf{l}_i + b)$ and
$C=10$. The objective is convex and solved by LBFGS. We then re-fit
every base learner on the \emph{full} training set (with $\mathbf{w}$)
to obtain the deployment models and collect their OOF-mode meta outputs
$\ptilde_{\mathrm{oof},i}$. The isotonic calibrator $g$ is fit on
$(\ptilde_{\mathrm{oof},i}, y_i)$ via the Pool-Adjacent-Violators
algorithm to enforce monotonicity. Finally, the decision threshold is
selected as
\[
\taustar \;=\; \argmax_{\tau\in(0,1)}\;
\mathrm{BalAcc}\!\bigl(\{\phat_{\mathrm{oof},i}\}, y\,;\tau\bigr),
\qquad \taustar = 0.135,
\]
where $\phat_{\mathrm{oof},i} = g(\ptilde_{\mathrm{oof},i})$. All three
quantities $g$, $\taustar$, and the bundle $\{\hat{\bm\beta},\hat b\}$
are \emph{frozen} after this step and never re-estimated on test data.

\paragraph{Stage 4: Deployment artifact bundle.}
The final exported objects are: (a) the five refit base learners
$\{f_k\}_{k=1}^{5}$, (b) the meta coefficient tuple $(\hat{\bm\beta},\hat b)$,
(c) the isotonic map $g$, and (d) the threshold $\taustar$. These are
precisely the four categories of objects consumed at inference
(Figure~\ref{fig:inference}).

\begin{figure}[t]
  \centering
  \includegraphics[width=\linewidth]{fig_B_training.pdf}
  \caption{\textbf{Training protocol of the champion model.}
  \textbf{Stage 1 (Cleanlab weighting).} A noise-robust probe classifier
  computes per-sample self-confidence on an independent CV loop; the
  bottom 10\% of samples are assigned $w_i=0.3$, the remainder
  $w_i=1.0$, producing the binary weight vector $\mathbf{w}$.
  \textbf{Stage 2 (5-fold OOF).} Each of the five base learners is
  trained under stratified 5-fold CV with $\mathbf{w}$ injected; the
  concatenated held-out predictions form the leakage-free OOF matrix
  $\Pof\in[0,1]^{N\times 5}$.
  \textbf{Stage 3 (Meta fit + Isotonic).} $\Pof$ is logit-transformed
  and passed to an L2-regularized logistic regression that minimizes
  Equation~\ref{eq:meta-objective} (yellow box, solved by LBFGS) to
  recover $(\hat{\bm\beta},\hat b)$; each base learner is then refit on
  the full training set; an isotonic calibrator $g$ is fit on
  $(\ptilde_{\mathrm{oof}}, y)$ via Pool-Adjacent-Violators; the
  decision threshold $\taustar$ is selected by maximizing balanced
  accuracy on the OOF predictions.
  \textbf{Stage 4 (Frozen artifacts).} The exported bundle
  $\{f_k,\hat{\bm\beta},\hat b,g,\taustar\}$ is exactly the set of
  objects consumed at inference time in Figure~\ref{fig:inference}.}
  \label{fig:training}
\end{figure}

% ------------------------------------------------------------
\subsection{Bridging Figures~\ref{fig:inference} and~\ref{fig:training}}
\label{sec:methods-bridge}

Figures~\ref{fig:inference} and~\ref{fig:training} depict two
complementary views of the same model system. Figure~\ref{fig:inference}
fixes all parameters and shows how an input traverses the frozen
pipeline to produce a binary alert; Figure~\ref{fig:training} reveals
how all learnable components -- the five base learners' internal
weights, the meta coefficients $(\bm\beta,b)$, the isotonic map $g$, and
the threshold $\taustar$ -- are jointly estimated from training data.
The two figures share (i)~the same 41-dimensional feature space
$\varphi$, (ii)~the same sample-weight vector $\mathbf{w}$,
(iii)~the same base-learner hyperparameters, and (iv)~the same 5-fold
split. Every artifact produced by Stage~4 of Figure~\ref{fig:training}
has exactly one consumption site in Figure~\ref{fig:inference}:
$(\bm\beta,b)$ is consumed only at the meta-$\sigma$ node, $g$ only at
the isotonic node, and $\taustar$ only at the threshold node. This
bijective training--inference correspondence implies that model
reproducibility reduces to the integrity of a single artifact bundle
hash, and that the OOF construction in Stage~2 guarantees the
distribution of base probabilities seen by the meta-learner at training
time coincides with that seen at deployment time.

% ------------------------------------------------------------
\subsection{Hard-Constrained Counterfactual (HRCF) Intervention}
\label{sec:methods-hrcf}
% TODO — use docs/paper/06_methods_hrcf.md as source
\textit{Placeholder. Projection-based counterfactual generation under
hard actionability constraints; cost function; top-$k$ menu
enumeration; comparison with soft-constrained DiCE variants.}

\end{document}
